\section{Introduction}

\subsection{Purpose}
This functional specification defines the functional and performance
requirements for the HVDC transmission system to be supplied for the
\emph{Utah/Nevada Industrial Grid (UNIG)} scheme and is the basis for HVDC system supplier selection.

\subsection{Scope}
\label{sec:scope}
The scope of supply is the HVDC terminals, together with any additional
equipment that the supplier considers necessary for the integration of the
HVDC system into the grid.

The DC line route and the DC tower design are under development and are not
fixed at the time of issue of this specification. The resulting DC line
parameters will be provided as they become available; for the associated dates
refer to Section~\ref{sec:timeline}.

\subsection{Key project characteristics and design drivers}
The following characteristics drive the design and execution of the scheme and
should be considered by the supplier:
\begin{itemize}\setlength{\itemsep}{0.35em}
  \item \textbf{Islanded operation.} The system operates as an islanded network
        with \emph{no interconnection to the wider AC grid}. The HVDC link and
        its connected generation and loads form an islanded network,
        which the converter control must be able to establish and stabilise.
  \item \textbf{Staged generation and transmission buildout.} Generation is
        built out in 300\,MW steps, with the pace set by water
        availability. staged buildout of the transmission line, phased
        to match the generation ramp, is therefore beneficial and is a design
        objective.
  \item \textbf{High reliability requirements.} The data-center loads impose
        high reliability and availability requirements on the
        transmission system (see Section~\ref{sec:availability}).
  \item \textbf{Standardized hardware design.} A standardized hardware
        design approach shall be adopted to reduce hardware lead times and to
        limit configuration variations between schemes.
\end{itemize}



